\subsection{\textbf{Sprint 0}}

O projeto teve início com uma reunião com as \textit{stakeholders} Eliza Pasquali Lunardi e Letícia Fossá Rosa, onde tivemos a oportunidade de entender com mais profundidade do que se tratava o Estudify. Nesse encontro, tiramos dúvidas de escopo e prioridades, e alguns ajustes importantes emergiram em relação ao material de abertura do projeto: a revisão de redação, que constava como item do escopo, acabou ficando de fora, enquanto a página web de administração, inicialmente listada como fora do escopo, foi incorporada. Saímos da reunião com bastante material para trabalhar ao longo da \textit{sprint}.

Logo em seguida, nos reunimos entre os \ac{ages} IV para alinhar como conduziríamos o projeto. Já no primeiro dia, debatemos alguns princípios que nortearam toda a nossa atuação. Definimos que as \textit{tasks} seriam escritas de forma bem mais detalhada do que o padrão --- uma decisão baseada no nosso histórico na \ac{ages}, que mostrava que tarefas mal definidas designadas a \ac{ages} I ou II tendiam a gerar retrabalho significativo. Também definimos que trabalharíamos organizados em \textit{squads}, compostas por pelo menos um integrante de cada senioridade, e que cada \ac{ages} IV ficaria responsável por acompanhar uma delas, com a missão de garantir que o time avançasse nas entregas. Essa divisão nos deu um controle muito mais assertivo sobre o andamento do projeto.

Para fundamentar a escolha do \textit{stack}, elaboramos uma pesquisa interna mapeando as tecnologias mais familiares ao time. Deixamos claro desde o início que a decisão final caberia aos \ac{ages} III, por serem eles os responsáveis pela arquitetura. Após a discussão técnica, definiu-se React Native para o \textit{mobile} e NestJS para o \textit{backend}; em conversa com os \ac{ages} II, que responderiam pelo banco, fechamos em PostgreSQL.

Quando conversamos com os \ac{ages} III sobre a distribuição de responsabilidades na \textit{sprint}, dois deles, Eduardo Camana e Isadora Guerra, manifestaram interesse em se dedicar inteiramente à prototipação. Ambos vinham do Apple Academy, e como eu já havia tido experiências muito positivas com a prototipação de quem participa deste programa, decidimos aceitar. Ainda assim, abrimos o convite para qualquer outro integrante do time que quisesse aprender mais sobre essa frente se juntar a eles. Os três \ac{ages} III restantes assumiram a montagem da infraestrutura. Esse tipo de paralelismo só foi viável porque tínhamos cinco \ac{ages} III, um número acima da média para projetos da \ac{ages}, e uma equipe muito engajada no geral.

Em paralelo, estruturamos o escopo do aplicativo em nove épicos, cada um agrupando um conjunto coeso de \ac{us}. Fiquei responsável por escrever as \ac{us} de dois deles: o Épico 3, que concentrava o \textit{core} do produto, toda a jornada de resolução de questões, e o Épico 6, referente à gamificação do aplicativo como um todo. Reforçamos entre nós a disciplina de escrever \textit{tasks} o mais detalhadas possível, com critérios de aceite claros e contratos bem definidos entre \textit{front} e \textit{back}. Quando terminamos de redigir as histórias, nos reunimos como \ac{ages} IV e revisamos uma a uma, garantindo que cada ponto mapeado no escopo tivesse ao menos uma \textit{task} correspondente no ClickUp.

Um ponto que não previmos foi o ritmo de entrega dos \ac{ages} II no modelo de banco. Nossa expectativa era de que um modelo minimamente concreto só apareceria por volta do meio da \textit{Sprint} 1, e estruturamos a escrita das \ac{us} partindo dessa premissa. Os \ac{ages} II, no entanto, conseguiram entregar um modelo bastante completo já no fim da \textit{Sprint} 0, o que, apesar de ser um ótimo resultado de engajamento da equipe, fez com que algumas das histórias que escrevemos ficassem ligeiramente desalinhadas em relação ao modelo que chegou. Não foi um erro muito significativo, mas gerou um ruído extra no início da \textit{Sprint} 1 que poderia ter sido evitado. Ficou como aprendizado, para mim como \ac{ages} IV, a necessidade de sincronizar mais de perto a escrita das \ac{us} com o avanço da modelagem de dados, inclusive nos cenários em que o time técnico entrega acima do esperado, o descompasso pode vir pelos dois lados.

Na apresentação de fechamento da \textit{sprint}, as \textit{stakeholders} ficaram bastante satisfeitas com os protótipos no geral; o resultado foi muito além do que haviam imaginado. Solicitaram apenas alguns ajustes pontuais na \ac{ui}, como trocar a fonte por uma com traços mais arredondados e deixar as telas mais coloridas. Ficaram também de nos enviar alguns ícones que gostariam de ver aplicados no aplicativo, bem como uma lista de questões já existentes no banco de dados delas, que serviria de base para o conteúdo inicial.

Ao final da \textit{sprint}, senti que o time em geral fez um excelente trabalho. O único ponto de atenção que levantei junto aos outros \ac{ages} IV foi que os \ac{ages} I acabaram tendo pouca participação nessa \textit{sprint}: como normalmente são eles que contribuem com os \textit{mockups} nessa fase, e essa frente ficou a cargo do Eduardo e da Isadora, sobrou menos espaço de atuação para eles. Não foi um problema de engajamento, mas algo a observar nas próximas \textit{sprints} para garantir que essa senioridade tenha frentes claras em que possa contribuir.
